% =============================================================================
%  Sección 3: Conclusiones
%  Sustituye el contenido de \lipsum con tu propio texto.
% =============================================================================

\section{Conclusiones}

En este ensayo hemos analizado varias perspectivas y conceptos que combinan desde lo 
meramente teórico en sustentabilidad y lo matemático en indicadores de análisis, 
en su mayoría hechos con datos recolectados de el portal de datos abiertos y las tecnologías 
implementados como Apimetro y VFTModel. Por lo que dividiremos las conclusiones en tres sectores:

\begin{itemize}
    \item \textbf{Sustentable-Social:} Concentrándonos mayormente en el como el transporte y su 
        robustez (o falta de ella) afecta al tejido social en términos actuales y de largo plazo.            
    \item \textbf{Sustentable-Urbano:} Tomando los conceptos del Modelo DOT, concluiremos si estos 
        se cumplen en el estado actual del transporte de la \zmvm y como la implementación de un o 
        unos corredores anillares, beneficiarían a dicha área.
    \item \textbf{Sustentable-Gobierno:} Tomando el punto anterior, definiremos lo faltante y lo que 
        se ha trabajado correctamente en el estado actual de la movilidad y como una integración 
        de gobierno federal sería lo más viable.
\end{itemize} 
\subsection{Sustentable a nivel Social}

Bajo esta perspectiva, la movilidad no solo representa el transporte de grandes cantidades de personas de 
punto A a punto B; si no un espacio intermedio entre el hogar de la persona y el lugar donde desarrolla sus 
actividades del día a día. Esto no solo afecta significativamente a las personas que viven en zonas periféricas 
o que sus condiciones laborales o educativas no se encuentran en los \termino{CBD} conectados en las alcaldías 
centrales de la \zmvm. Revela una vulnerabilidad de la población para llegar a sus destinos 
y una pérdida de la eficiencia en cuestiones de tiempo y seguridad. \\ Como toda metrópoli, la 
\zmvm se expande a cada año que pasa y la cantidad de personas que residían en 1960 no corresponde a las mismas 
que llegaron en 2010 y mucho menos a la fecha de escrito este ensayo (2026). Prever un control de natalidad de la 
población y sus asentamientos irregulares es solo un síntoma de la dispersión de vivienda de la población 
circundante que llega a establecerse para trabajar en la metrópoli mexicana.

\subsection{Sustentable a nivel Urbano}

Como se exploró en \text{\termino{Sustentable a Nivel Social}}, el crecimiento de la población genera un fenómeno 
conocido como \termino{Leap Frog} o \termino{Salto de Rana} \cite{jia2010}. Fenómeno que provoca un crecimiento 
no uniforme y desordenado. Esto se convierte en un problema a nivel logístico, como proveer de transporte seguro, 
eficiente y accesible a la población naciente de estos lugares fragmentados, y más aún, como poder calcular el impacto 
que este tendrá en un futuro usando los términos de sostenibilidad en el urbanismo \termino{Económico, Ecológico y Social}.
\\ La respuesta es difícil encontrarla en informes oficiales o en documentos de países con condiciones diferentes. 
México es uno de los grandes de Latinoamérica, pero su condición geográfica y económica, al igual que su historia y 
población, vuelven muy única (y difícil) la implementación de una solución de escritorio.

\subsection{Sustentable a nivel Gobierno}

A nivel geográfico, el Estado de México, es casi dos veces mayor que la Ciudad de México, sin embargo, en términos 
políticos y de vivienda, la Ciudad de México es casi 3 veces más grande. Siendo esta la entidad que recibe más presupuesto 
federal anual y por sexenio. Esto se debe principalmente a la centralización de los tres poderes que rigen a la Nación. 
A pesar de elo, la inferencia de la Ciudad de México no se queda en su entidad geográfica, ni siquiera en su vecino cercano 
(Estado de México), se reparte en los estados circundantes (Morelos, Tlaxcala y parte de Pachuca), lo que revela que la 
administración del gobierno no puede quedarse solamente en su entidad. Es cierto que existen consejos, comunicación y una mutua 
cooperación de los mandatarios por entenderse y organizarse, pero esto en su mayoría depende del partido del que provengan y 
las rivalidades (o trivialidades) que puedan llegar a tener. Básicamente es como una moneda al aire, cada 6 años (en promedio) 
experimentamos un cambio de gobierno y con él, un \termino{reset} de las relaciones públicas de la \zmvm. \\
La propuesta más acertada a este problema es la creación de un gobierno con injerencia metropolitana, o bien una secretaria 
con presupuesto federal que se encarga de administrar y designar funcionarios capaces para su administración. Mientras 
esta solución solo sea un imaginario, los problemas de toda la zona seguirán encerrados en feudos regionales llamados 
\termino{estados federales}.